\documentclass[a4paper,11pt]{article}
\usepackage[utf8]{inputenc}
\usepackage[T1]{fontenc}
\usepackage[english]{babel}
\usepackage[left=1cm,right=1cm,top=.5cm,bottom=0cm]{geometry}
\usepackage{enumitem}
\usepackage{hyperref}
\usepackage{xcolor}
\usepackage{titlesec}
\usepackage{marvosym}
\usepackage{lmodern}
\usepackage[normalem]{ulem}

\definecolor{primary}{RGB}{0, 51, 102}
\definecolor{secondary}{RGB}{80, 80, 80}

\hypersetup{colorlinks=true, linkcolor=primary, filecolor=primary, urlcolor=primary}
\titleformat{\section}{\large\bfseries\color{primary}\uppercase}{}{0em}{}[\titlerule]
\titlespacing{\section}{0pt}{8pt}{5pt}

\begin{document}
\begin{center}
    {\Large \textbf{Lo\"ic Aron MBASSI EWOLO}} \\ \vspace{4pt}
    \textbf{\large Alternant Ing\'enieur DevOps -- Cloud \& Infrastructure} \\
    \textit{\small Stage INRIA Saclay (Avril--Ao\^ut 2026) -- Disponible en alternance d\`es septembre 2026 (12 mois)} \\ \vspace{4pt}
    \small
    \Pointinghand\ Cergy, France (Mobile IDF) \quad \Mobilefone\ (+33) 7 59 52 33 98 \quad
    {\color{primary}\Letter\ \href{mailto:loicaron.mbassiewolo@ensea.fr}{loicaron.mbassiewolo@ensea.fr}} \\ \vspace{2pt}
    {\color{primary}LinkedIn: \href{https://www.linkedin.com/in/loic-mbassi/}{linkedin.com/in/loic-mbassi}} \quad
    {\color{primary}GitHub: \href{https://github.com/Nameless0l}{github.com/Nameless0l}} \quad
    {\color{primary}Portfolio: \href{https://mbassiloic.tech}{mbassiloic.tech}}
\end{center}
\vspace{-6pt}
\noindent\small{Mise en place de pipelines CI/CD et containerisation Docker dans le cadre du projet PRIMO \`a l'INRIA Saclay, et architecture r\'eactive microservices (Spring WebFlux, R2DBC, Redis) pour la plateforme Snappy supportant des milliers de connexions simultan\'ees. Formation double dipl\^ome ENSEA/Polytechnique Yaound\'e en syst\`emes et informatique avanc\'ee. Exp\'erience en orchestration de services (RabbitMQ, API Gateway, Docker Compose) dans le projet GoFind. TF1, g\'eant de l'audiovisuel, exige une infrastructure DevOps fiable pour la diffusion de contenus \`a grande \'echelle.}

\section{Formation}
\begin{itemize}[leftmargin=*, label=\tiny$\bullet$, nosep]
    \item \textbf{ENSEA} \hfill \textit{2025 -- 2027} \\
    \small{Double dipl\^ome d'Ing\'enieur -- Syst\`emes, Informatique Avanc\'ee \& DevOps.}
    \item \textbf{\'Ecole Polytechnique de Yaound\'e (1\textsuperscript{\`ere} \'ecole d'ing\'enieur CEMAC)} \hfill \textit{2021 -- 2025} \\
    \small{Dipl\^ome d'Ing\'enieur de Conception (Master 2) : \textbf{Syst\`emes d'Information}, Architecture Logicielle, G\'enie Logiciel.}
\end{itemize}

\section{Comp\'etences Techniques}
\begin{itemize}[leftmargin=*, label=\tiny$\bullet$, nosep]
    \item \textbf{DevOps/CI-CD :} Docker, Docker Compose, Git, pipelines CI/CD (GitHub Actions, Jenkins concepts), automatisation builds/d\'eploiements, Makefile, versioning.
    \item \textbf{Cloud/Infrastructure :} Containerisation, microservices, orchestration (Kubernetes concepts), monitoring, load balancing, API Gateway, haute disponibilit\'e.
    \item \textbf{Backend/APIs :} Spring Boot 3, Spring WebFlux, Laravel, Node.js, REST APIs, gestion transactions, optimisation requêtes.
    \item \textbf{Bases de Donn\'ees :} PostgreSQL, MySQL, Redis, MongoDB, R2DBC, mod\'elisation, indexation, cache distribu\'e.
    \item \textbf{Langues :} Fran\c{c}ais natif, Anglais professionnel (INRIA, documentation technique EN).
\end{itemize}

\section{Exp\'eriences \& Projets}
\begin{itemize}[leftmargin=*, label=\tiny$\bullet$, nosep]
    \item \textbf{Stage INRIA Saclay -- \'Equipe TRIBE (Projet PRIMO)} \hfill \textit{Avril-Ao\^ut 2026} \\
    \textit{Palaiseau (IP Paris)} \hfill \textit{Python, Docker, Git, CI/CD, pipelines de d\'eploiement, NLP}
    \vspace{-2pt}
    \begin{itemize}[leftmargin=0.4cm, nosep, label=\tiny$\bullet$]
        \item \small{Mise en place pipeline Python reproducible : containerisation Docker (Dockerfile, Docker Compose), CI/CD avec versioning Git.}
        \item \small{Orchestration services (BERT, scikit-learn, pandas) ; automatisation testing, d\'eploiement predictable, monitoring logs.}
        \item \small{Collaboration IP Paris/LVMT ; documentation technique EN, code infrastructure-as-code.}
        \item \small{Exp\'erience fondatrice en DevOps ML : packaging mod\`eles, versioning, d\'eploiement production-ready.}
    \end{itemize}
    \vspace{3pt}

    \item \textbf{Snappy -- Plateforme Messagerie Enterprise (Architecture DDD)} \hfill \textit{2024-2025} \\
    \textit{ENSEA} \hfill \textit{Spring Boot 3, Spring WebFlux, PostgreSQL, R2DBC, Redis, Docker, Socket.IO}
    \vspace{-2pt}
    \begin{itemize}[leftmargin=0.4cm, nosep, label=\tiny$\bullet$]
        \item \small{Architecture r\'eactive Spring WebFlux : gestion milliers connexions WebSocket non-blocking, backpressure, asynchrone.}
        \item \small{Persistance haute-perf : R2DBC (PostgreSQL driver r\'eactif), Redis cache distribu\'e, transactions concurrentes ACID.}
        \item \small{Docker containerisation, CI/CD ready, orchestration compose pour environnements multi-services.}
        \item \small{DDD (Domain-Driven Design) : structuration code modulaire, s\'eparation concerns backend/infrastructure.}
    \end{itemize}
    \vspace{3pt}

    \item \textbf{Dev Backend CSC -- Fintech/Bancaire} \hfill \textit{Oct 2024-Jan 2025} \\
    \textit{Yaound\'e} \hfill \textit{Laravel, MySQL, REST APIs, Docker, Redis, Git, CI/CD, tests}
    \vspace{-2pt}
    \begin{itemize}[leftmargin=0.4cm, nosep, label=\tiny$\bullet$]
        \item \small{Architecture microservices transactions : gestion concurrence (locks pessimistes, versioning optimiste), scalabilit\'e.}
        \item \small{APIs RESTful document\'ees (Swagger), CI/CD Git-native (hooks, automated testing), d\'eploiement reproductible.}
        \item \small{Caching Redis pour haute disponibilit\'e, indexation MySQL pour perf requ\^etes, load testing.}
    \end{itemize}
    \vspace{3pt}

    \item \textbf{GoFind -- Architecture Microservices Orchestr\'ee} \hfill \textit{2024} \\
    \textit{ENSEA} \hfill \textit{Laravel, NestJS, MongoDB, RabbitMQ, Docker, API Gateway, Docker Compose}
    \vspace{-2pt}
    \begin{itemize}[leftmargin=0.4cm, nosep, label=\tiny$\bullet$]
        \item \small{Architecture microservices : Laravel backend, NestJS service, MongoDB, orchestration RabbitMQ async.}
        \item \small{API Gateway pour routing, CQRS patterns, Event Sourcing ; orchestration services via Docker Compose.}
        \item \small{Scalabilit\'e horizontale, service discovery, r\'esilience (circuit breakers, retries, timeouts).}
    \end{itemize}
    \vspace{3pt}

    \item \textbf{Urban Waste Management -- Smart City (1\textsuperscript{er} Prix National)} \hfill \textit{2024} \\
    \textit{ENSEA} \hfill \textit{Python, Docker, pandas, scikit-learn, d\'eploiement}
    \vspace{-2pt}
    \begin{itemize}[leftmargin=0.4cm, nosep, label=\tiny$\bullet$]
        \item \small{Pipeline ML end-to-end : containerisation Docker, mod\`eles pr\'edictifs IoT, d\'eploiement reproductible.}
        \item \small{1\textsuperscript{er} prix CITS National Hackathon (vs 75 \'equipes) ; infrastructure production-ready.}
    \end{itemize}
    \vspace{3pt}

    \item \textbf{Assistant de Recherche @ MILA} \hfill \textit{Juin-Sept 2025} \\
    \textit{Remote (Universit\'e Montr\'eal)} \hfill \textit{Python, PyTorch, TensorFlow, notebooks Jupyter, versioning}
    \vspace{-2pt}
    \begin{itemize}[leftmargin=0.4cm, nosep, label=\tiny$\bullet$]
        \item \small{Exp\'eriences ML reproductibles : versioning code, seeds fixés, documentation, notebooks annotés.}
        \item \small{Exp\'erience hands-on en MLOps : tracking experiments, visualization, reporting.}
    \end{itemize}
\end{itemize}

\section{Distinctions \& Leadership}
\begin{itemize}[leftmargin=*, label=\tiny$\bullet$, nosep]
    \item \small{\textbf{1\textsuperscript{er} prix CITS National Hackathon} (Smart City, infrastructure ML) -- vs 75 \'equipes.}
    \item \small{\textbf{Head of Projects Cell \& Lead AI Cell ENSEA} -- +50\% membres actifs, architecture projets.}
    \item \small{\textbf{Open Source :} Laravel API Generator (+30 \'etoiles GitHub, Packagist) -- infrastructure code generation.}
    \item \small{\textbf{Certifications :} Supervised ML Stanford/DeepLearning.AI (Oct. 2024), Python for Data Science IBM (Jan. 2024).}
\end{itemize}

\end{document}
